\documentclass[10pt,a4paper]{amsart}
\usepackage[margin=19mm]{geometry}
\usepackage{amsmath,amssymb,mathtools}
\usepackage[scr=rsfs,cal=euler]{mathalfa}
\usepackage{xfrac}
\usepackage[unicode,hidelinks,bookmarks=false]{hyperref}
\newcommand{\ie}{i.e.\ }
\newcommand{\cf}{cf.\ }
\newcommand{\resp}{respectively\ }
\newcommand{\Subsection}{\subsection}
\providecommand{\nobreakdash}{\nobreak\hbox{-}}
\setlength{\emergencystretch}{2em}
\multlinegap0pt
\usepackage{bm}
\usepackage{bbm}
\usepackage{dsfont}
\usepackage{xspace}
\usepackage{cancel}
\usepackage{mathtools}
\usepackage{amsthm}
\makeatletter
\@ifundefined{theorem}{\newtheorem{theorem}{Theorem}}{}
\@ifundefined{lemma}{\newtheorem{lemma}[theorem]{Lemma}}{}
\makeatother
\newcommand{\C}{{\mathbbm C}}
\newcommand{\N}{{\mathbbm N}}
\newcommand{\Q}{{\mathbbm Q}}
\newcommand{\Z}{{\mathbbm Z}}
\renewcommand{\eqref}[1]{(\ref{#1})}   
\title[Irreducibility of eventually $2$-periodic curves in the modu]{Irreducibility of eventually $2$-periodic curves in the moduli space of cubic polynomials}
\author{Niladri Patra}
\date{Source: arXiv:2305.19944v2 (CC BY 4.0)}
\begin{document}
\maketitle

\noindent\emph{Source and licence.} Irreducibility of eventually $2$-periodic curves in the moduli space of cubic polynomials. Version arXiv:2305.19944v2 (CC BY 4.0), distributed under the Creative Commons Attribution 4.0 International licence, as recorded in the arXiv licence metadata for this exact version and shown on the abstract page https://arxiv.org/abs/2305.19944v2. Full rights notice: licenses/arxiv-2305.19944-CC-BY-4.0.txt.\par\smallskip

\noindent\textit{Editorial scope.} Counted unit: the Lemma whose source label is \texttt{k23res} in the article (source0.tex lines 431--438, source label \texttt{k23res}), delivered together with the article's own complete original proof of it (source0.tex lines 439--507, one \texttt{proof} environment of 69 lines). No mathematics is written, repaired or re-derived: statement and proof are the article's own text line for line. Required context retained uncounted: k (lines 244--246); thurs (lines 255--261); maineqcopy (lines 346--348); h123 (lines 356--364); h023. (lines 367--381); h023 (lines 391--393); divlemma (lines 415--417). Every definition, display and label of those lines is retained unchanged. Omitted from this package and not redistributed: 1--243, 247--254, 262--345, 349--355, 365--366, 382--390, 394--414, 418--430, 508--714. Nothing inside a retained excerpt is omitted or altered: every retained body is delivered exactly as the article's lines are written, so no omission receipt of any kind accompanies this package. Citations: the retained excerpts carry no citation command at all, so no bibliography entry of the article is transcribed and no reference is redistributed. Reference adaptation: none was needed and none was made; every reference inside the delivered unit resolves inside it, so no unresolved reference or citation is produced. Printed slips kept literal: no printed word, symbol, hypothesis or number of the counted statement or of its proof was altered. Layout and class: the article's own class is \texttt{amsart} with packages including amsmath, amssymb, latexsym, amsthm, amsfonts, mathrsfs, amscd, color, bbm, scalerel, stackengine, microtype, tikz-cd, setspace; the wrapper is the lane's pinned \texttt{amsart} editorial wrapper at 10pt on A4, which restates the theorem-like environments the retained text needs and adds the article's own macro definitions that the retained text needs (\texttt{\textbackslash C}, \texttt{\textbackslash N}, \texttt{\textbackslash Q}, \texttt{\textbackslash Z}, \texttt{\textbackslash eqref}), with extra wrapper packages (bm, bbm, dsfont, xspace, cancel, mathtools). The delivered numbering is the unit's own. Assets: no retained span contains a figure environment, a graphics-inclusion command, a PDF-page inclusion, a table or a TikZ picture; no image file is copied into this package, the package declares no asset dependency and no image file is redistributed with it. Licence: version arXiv:2305.19944v2 of this article is distributed under the Creative Commons Attribution 4.0 International licence, as recorded in the arXiv licence metadata for this exact version and shown on the abstract page https://arxiv.org/abs/2305.19944v2. A full rights notice is delivered at licenses/arxiv-2305.19944-CC-BY-4.0.txt. Characters: every retained excerpt is pure ASCII, so the delivered engine, XeLaTeX with its default Latin Modern text font, reports no missing character.
    \begin{lemma}\label{k}
        Let $m,n\in \N$, and g.c.d.$(m,n)=r$. Let $g$ be an irreducible element of $\Z[a,b]$, monic as a polynomial in  $\Z[a][b]$. If $g$ divides both $f_{k,n}$ and $f_{l,m}$ in $\Z[a,b]$, then $g$ divides $f_{l,r}$ in $\Z[a,b]$.
    \end{lemma}

    \begin{theorem}\label{thurs}
        Fix $k_{1},k_{2}\in \N \cup \{0\}$, and $n_{1},n_{2}\in \N$. Then, the polynomials 
        \begin{equation*}
            f^{k_{1}+n_{1}}(a) - f^{k_{1}}(a) \text{ and } f^{k_{2}+n_{2}}(-a) - f^{k_{2}}(-a)
        \end{equation*}
        are coprime in $\C[a,b]$.
    \end{theorem}

    \begin{equation}\label{maineqcopy}
        f(z) = z^{3} - 3a^{2}z + 2a^{3} + b, 
    \end{equation}

    \begin{lemma}\label{h123}
        The polynomial $h_{1,2}$ is $(b-a)^{2}+1$, which has the following properties:
        \begin{itemize}
            \item It is irreducible over $\Q$,
            \item It is reducible and smooth over $\C$,
            \item There is no $\Q$-rational point on the solution curve of $h_{1,2}$ in $\C^{2}$,
            \item The curve, $\mathscr{S}_{1,2}$, of $h_{1,2}$ in $\mathcal{M}_{3}$ is an irreducible line.
        \end{itemize}
    \end{lemma}

        \begin{align}
            f_{0,2} & =f(b)-a \notag\\
            & = b^{3} -3a^{2}b+2a^{3}+b-a \notag\\
            & = (b-a)\left((b-a)(b+2a)+1\right). \label{h023.}\\
            f_{1,1} & =f(b)-b \notag \\
            & =b^{3}-3a^{2}b+2a^{3}\notag \\
            & =(b-a)^{2}(b+2a).\label{f113}\\
            f_{1,2} & =f^{3}(a)-f(a)\notag \\
            & =(f^{2}(a)-a)\left((f^{2}(a))^{2}+af^{2}(a)-2a^{2}\right)\notag \\
            & = f_{0,2} \left((f(b))^{2}+af(b)-2a^{2}\right)\notag \\
            & = f_{0,2}\; (f(b)-a)(f(b)+2a)\notag \\
            & = (f_{0,2})^{2} \;(f(b)+2a)\label{f123byf023}\\
            & = (f_{0,2})^{2}\left(b^{3}-3a^{2}b+2a^{3}+b+2a \right)\notag \\
            & = (f_{0,2})^{2} (b+2a)\left((b-a)^{2}+1\right).\label{f123}
        \end{align}

    \begin{lemma}\label{h023}
        The polynomial $h_{0,2}$ is $(b-a)(b+2a)+1$, and it is irreducible over $\C$.
    \end{lemma}

    \begin{lemma}\label{divlemma}
        For any $k\in \N$, we have $o(f_{k,2},f_{0,2})\geq 2$. For any even $k\in \N$, we have $o(f_{0,k},f_{0,2})=o(f_{0,k},h_{0,2})=1$.
    \end{lemma}

    \begin{lemma}\label{k23res}
        Let $l=b-a+i \in \C[a,b]$. Up to multiplication by a power of $i$, The resultant $\mbox{Res}(h_{k,2},l)$ is,
        \begin{center}
            $\mbox{Res}(h_{k,2},l)= \left\{ \begin{array}{rcl}
                3(2ai+1); & k \mbox{ even,} & k>0, \\ 3a; & k \mbox{ odd,} & k>1. \\
            \end{array}\right.$
        \end{center}
    \end{lemma}

    \begin{proof}
        Let $k\in \N, k > 1$. We will first remove irreducible factors of $f_{k-1,2}$ in $\Z[a,b]$, from $f_{k,2}$ with each such factor raised to the highest power that divides $f_{k,2}$.
        Consider the polynomial, 
        \begin{align}\label{hk231}
             g_{k}(a,b):= \frac{f_{k,2}}{f_{k-1,2}} & = \frac{f^{k+1}(b)-f^{k-1}(b)}{f^{k}(b)-f^{k-2}(b)}\notag \\
             & = \left(f^{k}(b)\right)^{2}+f^{k}(b)f^{k-2}(b)+\left(f^{k-2}(b)\right)^{2}- 3a^{2}\notag \\
             & \equiv 3\left(\left(f^{k-2}(b)\right)^{2}-a^{2}\right)\; \left( \text{mod } f_{k-1,2}\right).
        \end{align}
        So, any irreducible polynomial $s_{k}\in \Z[a,b]$ that divides both $g_{k}$ and $f_{k-1,2}$, will also divide 
        \begin{equation}\label{pereq}
            3\left((f^{k-2}(b))^{2}-a^{2}\right)=3 f_{0,k-1} (f^{k-2}(b)+a).
        \end{equation}
        From Thurston's rigidity theorem (Theorem \ref{thurs}), we get that $f^{k-2}(b)+a$  and $f_{k-1,2}$ are coprime. So, $s_{k}$ will divide $f_{0,k-1}$ (we can remove $3$, because $f_{k-1,2}$ is monic in $\Z[a][b]$, and so are its irreducible factors). As $s_{k}$ divides both $f_{0,k-1}$ and $f_{k-1,2}$, by Lemma \ref{k} we have that $s_{k}$ divides $f_{0,1}$ if $k$ is even, and $f_{0,2}$ if $k$ is odd.

        \textbf{Let $k\in \N$ be even.} As $f_{0,1}= b-a$, we get that 
        $$h_{k,2}\; \text{ divides }\; \frac{g_{k}(a,b)}{(b-a)^{i_{k}}},$$ where $i_{k}$ is the highest power of $(b-a)$ that divides $g_{k}$. Also, 
        $$\frac{g_{k}(a,b)}{(b-a)^{i_{k}}}\;\text{ is coprime to }\;f_{k-1,2}.$$

        \textbf{Let $k\in \N, k>1, k$ odd.} 
        From Equation \eqref{h023.} and Lemma \ref{h023}, we know that, over $\Q$, the irreducible factors of $f_{0,2}$ are $h_{0,2}$ and $(b-a)$. From Equations \eqref{hk231}, \eqref{pereq} and Lemma \ref{divlemma}, we have $o(g_{k},h_{0,2})=1$, for all $k\in \N, k>1$. So, 
        $$h_{k,2}\; \text{ divides }\; \frac{g_{k}(a,b)}{h_{0,2}\cdot (b-a)^{i_{k}}},$$ 
        where $i_{k}$ is the highest power of $(b-a)$ that divides $g_{k}$. Also,
        $$\frac{g_{k}(a,b)}{h_{0,2}\cdot (b-a)^{i_{k}}}\;\text{ is coprime to }\;f_{k-1,2}.$$

        For $k\in \N, k>1$, let us define \begin{equation}\label{hk232}
            g'_{k}(a,b) := \left\{ \begin{array}{rcl}
                g_{k}(a,b)/(b-a)^{i_{k}}; & k &\mbox{even} \\ g_{k}(a,b)/(h_{0,2}\cdot (b-a)^{i_{k}}); & k & \mbox{odd} \\
            \end{array}\right.
        \end{equation}

        Next, we will factor out the irreducible factors of $f_{k,1}$ from $g'_{k}(a,b)$. But irreducible factors of $f_{k-1,1}$ has already been factored out, since $f_{k-1,1}$ divides $f_{k-1,2}$. Hence, we need to consider the irreducible factors of 
        \begin{equation}\label{5.10}
            \frac{f_{k,1}}{f_{k-1,1}}=(f^{k}(a))^{2}+f^{k}(a)f^{k-1}(a)+(f^{k-1}(a))^{2}-3a^{2},
        \end{equation}
        and their highest powers that divide $g'_{k}(a,b)$. Let's denote the product of common irreducible factors of $f_{k,1}/f_{k-1,1}$ and $g'_{k}$, each raised to their highest power that divides $g'_{k}$, as $t_{k}(a,b)$. By definition of $h_{k,2}$,
        \begin{equation}\label{hk233}
            h_{k,2}(a,b) = \frac{g'_{k}(a,b)}{t_{k}(a,b)}.
        \end{equation}

        Now, we can compute the resultant. We have, $\mbox{Res}(h_{k,2},l)=h_{k,2}(a,a-i)$.

        Putting $b=a-i$ in Equation \eqref{maineqcopy}, by direct computation or by using Lemma \ref{h123} one obtains $f^{k}(a)=a-i, $ for odd $k\in \N$, and $f^{k}(a)=-2a, $ for even $k\in \N$. Moreover, $h_{0,2}(a,a-i)= -i(-i+3a)+1=-3ai$.

        Observe that for any $k\in \N$,
        \begin{align*}
            \left(\frac{f_{k,1}}{f_{k-1,1}}\right)(a,a-i) & = (a-i)^{2}-2a(a-i)+4a^{2}-3a^{2}\\
            & =-1.
        \end{align*}
        So, $t_{k}(a,a-i)=1$, up to multiplication by a power of $i$ (because, any irreducible factor $t'_{k}$ of $t_{k}$ in $\Z[a,b]$ divides $f_{k,1}/f_{k-1,1}$ in $\Z[a,b]$. So in $\Z[i][a]$, $t'_{k}(a,a-i)$ divides $(f_{k,1}/f_{k-1,1})(a,a-i)=-1$).

        From the above discussion, we get that, up to multiplication by a power of $i$,

        \textbf{For $k$ even,}
        \begin{align*}
            \mbox{Res}(h_{k,2},l) & =h_{k,2}(a,a-i) & \\
            & = g_{k}(a,a-i) & (\mbox{by Equation }\eqref{hk232})\\
            & = 3(a-i)^{2}-3a^{2} & (\mbox{by Equation }\eqref{hk231})\\
            & =-3(2ai+1). & 
        \end{align*}
        
        \textbf{For $k> 1, k$ odd,}
        \begin{align*}
            \mbox{Res}(h_{k,2},l) & = h_{k,2}(a,a-i) & \\
            & = \frac{g_{k}(a,a-i)}{h_{0,2}(a,a-i)} & (\mbox{by Equation }\eqref{hk232})\\
            & =\frac{3(4a^{2})-3a^{2}}{-3ai} & (\mbox{by Equation }\eqref{hk231})\\
            & = 3ai. & 
        \end{align*}
        Hence, the lemma is proved.
    \end{proof}

\end{document}
